% Cote 119: offprints, reports, courses, CV and letters (typed and printed part): archive cote 119, site page 180 (PDF page 181).
% typescript. Draft compared: none.
% AI-transcribed by Claude Opus 5.5 (claude-opus-5-5, adaptive thinking, effort high), Message Batches API, batch r119_full (msgbatch_01DwZNm5Y7YsKBWyAN4sFs8w), 2026-11-13; not yet checked by eye.
\scanpage{119/180}
\mspage{- 19 -}

est un chemin continu non nécessairement différentiable par morceaux.

6.8. \uline{Théorème}.

Si $\gamma_0$ et $\gamma_1$ sont deux chemins fermés dans U, homotopes comme
chemins fermés on a $\int_{\gamma_0} \omega = \int_{\gamma_1} \omega$ pour toute forme fermée $\omega$ dans U.

\uline{Démonstration} :

Soit $\delta$ vérifiant (\unsure{6}.2.1) et (6.2.3.) .

Soit f une primitive de $\omega$ suivant $\delta$. (théorème 6.5.). Pour u fixé
$t \to f(t,u)$ est une primitive de $\omega$ suivant $t \to \delta(t,u)$. On a donc (6.7)
$\int_{\gamma_0} \omega = f(1\unsure{,}0) - f(0,0)$
$\int_{\gamma_1} \omega = f(1,1) - f(0,1)$
Comme d'après (6.2.3) on a : $\delta(0,u) = \delta(1,u)$, $f(0,u)$ et $f(1,u)$ primitives
de $\omega$ suivant le chemin $u \to \delta(0,u)$ diffèrent d'une constante. En particulier
$f(1,1) - f(0,1) = f(1,0) - f(0,0)$ C.Q.F.D.

6.9. \uline{Corollaire}.

Toute forme différentielle fermée $\omega$ dans un ouvert simplement con-%
nexe U possède une primitive dans U.

\uline{Démonstration}.

D'après \unsure{(6,8)|(6.8)} on a $\int_{\gamma} \omega = 0$ pour tout chemin fermé $\gamma$. Car un
tel c hemin est homotope à un point. On conclut grâce à (3.3) .

\margin{180}
